\newpage
\begin{center}
  \textbf{1. Теоретические и правовые основы управления цифровыми активами граждан}
\end{center}
\refstepcounter{chapter}
\addcontentsline{toc}{chapter}{1. Теоретические и правовые основы управления цифровыми активами граждан}

\section{Понятие цифровых прав и активов, их классификация}

\subsection{Цифровые права и активы: определение и сущность}

Прежде чем рассматривать требования к системам управления цифровыми активами, необходимо четко обозначить область, в которой эти системы будут функционировать. В данном случае речь идет о цифровых правах и активах граждан.

Под цифровыми правами и активами понимаются права и активы, которые существуют в цифровой форме и могут быть использованы, переданы или управляемы с помощью цифровых технологий. К таким активам относятся, например, электронные документы, цифровые валюты, учетные записи в социальных сетях, цифровые фотографии и видео, а также интеллектуальная собственность в цифровом формате. [автор] определяет цифровые активы как "любые данные или информация, которые имеют ценность для владельца и существуют в цифровой форме"~\cite{Orlov2024Moluch}.

\subsection{Классификация цифровых активов}

Edwards Lilian предлагает рассматривать общий портфель цифровых активов, который является продолжением человеческого капитала. Мы не можем игнорировать, что в него входят не только цифровые финансовые активы, но и другие цифровые следы, которые человек оставляет в интернете: социальные сети, электронная почта, облачные хранилища и т.д.~\cite{Edwards2013PostMortemPrivacy}.

В работе рассматривается классификация цифровых активов, основанная на подходе предложенном в книге Э. Харбинья, расширенная с учётом добавления недостающих категорий, преимущественно цифровых финансовых активов~\cite{Harbinja2022DigitalDeath, Orlov2024Moluch}. в результате чего получилась следующая классификация цифровых активов:

\begin{landscape}
\begin{table}[htbp]
\centering
\small
\renewcommand{\arraystretch}{1.2}

\begin{tabular}{|p{2.3cm}|p{3.0cm}|p{2.2cm}|p{2.2cm}|p{3.2cm}|p{3.5cm}|}
\hline
\textbf{Класс актива} &
\textbf{Примеры} &
\textbf{Состав данных} &
\textbf{Ценность} &
\textbf{Правовая природа} &
\textbf{Посмертные операции и риски} \\
\hline

Коммуникации (email) &
письма, вложения, папки &
контент, метаданные &
личная, доказательная &
персональные данные, договор &
выборочная передача, архивирование; риск раскрытия данных третьих лиц \\
\hline

Сообщения и мессенджеры &
личные сообщения, чаты &
контент, метаданные &
личная &
персональные данные высокой чувствительности &
как правило недоступны наследникам; требуется прямое правовое основание \\
\hline

Социальные сети (публичный контент) &
посты, фото, видео, комментарии &
контент, метаданные &
мемориальная, репутационная &
IP, персональные данные, договор &
мемориализация, удаление; конфликты интересов родственников и платформ \\
\hline

Цифровой профиль &
аккаунт, bio, аватар &
учётная запись, контент &
идентичность &
персональные данные, договор &
изменение статуса, закрытие; платформа выступает фактическим арбитром \\
\hline

Социальный граф &
друзья, контакты, реакции &
метаданные &
реляционная &
персональные данные третьих лиц &
ограниченный экспорт; споры о допустимости передачи \\
\hline

Личные архивы &
заметки, фотоальбомы &
контент &
личная &
персональные данные &
передача доверенному лицу; возможны чувствительные данные \\
\hline

История активности &
локации, поиски, просмотры &
метаданные &
очень чувствительная &
персональные данные &
как правило подлежит удалению; высокий репутационный риск \\
\hline

Игровые и виртуальные аккаунты &
персонажи, предметы &
учётная запись, контент &
личная, квази-финансовая &
договор &
обычно непередаваемы по правилам платформ \\
\hline

\end{tabular}

\caption{Классификация нефинансовых цифровых активов граждан}
\label{tab:digital-assets-nonfinancial}
\end{table}
\end{landscape}

\begin{table}[htbp]
\centering
\small
\renewcommand{\arraystretch}{1.2}

\begin{tabular}{|p{2.6cm}|p{3.0cm}|p{2.2cm}|p{2.4cm}|p{3.4cm}|}
\hline
\textbf{Класс актива} &
\textbf{Примеры} &
\textbf{Состав данных} &
\textbf{Правовая природа} &
\textbf{Посмертные операции и риски} \\
\hline

Банковские счета &
карты, вклады, интернет-банк &
учётные записи, операции &
договор, право требования &
наследование по банковской процедуре; доступ по паролю недопустим \\
\hline

Брокерские счета &
ценные бумаги, портфели &
учётные записи, отчёты &
договор, финансовые права &
блокировка до вступления в наследство; строгая проверка полномочий \\
\hline

Платёжные сервисы &
электронные кошельки &
учётные записи, балансы &
договор &
ограничения KYC/AML; возможен отказ провайдера \\
\hline

Криптоактивы (self-custody) &
приватные ключи, seed phrase &
учётные данные &
фактический контроль ключа &
потеря ключа необратима; высокий риск компрометации \\
\hline

Криптоактивы (custodial) &
биржевые аккаунты &
учётные записи, балансы &
договор, регуляторные требования &
наследование через провайдера; юрисдикционные риски \\
\hline

Бонусы и баллы &
программы лояльности &
учётные записи &
договор &
часто непередаваемы; возможное сгорание \\
\hline

Монетизируемые аккаунты &
донаты, реклама, выплаты &
аккаунт, контент, финансы &
договор, IP, выплаты &
разделение прав на аккаунт, контент и доходы \\
\hline

\end{tabular}

\caption{Классификация финансовых цифровых активов граждан}
\label{tab:digital-assets-financial}
\end{table}

Кроме этого цифровые активы можно классифицировать по уровню доступа (личные, корпоративные, публичные), по способу хранения (локальные, облачные) и по степени ликвидности (высоколиквидные, низколиквидные).

Таким образом, понимание сущности и классификации цифровых прав и активов является фундаментальным для разработки эффективных систем управления этими активами. Необходимо учитывать разнообразие типов цифровых активов и их особенности, чтобы обеспечить надежное и безопасное управление ими.

\section{Правовые аспекты управления цифровыми активами граждан}

\subsection{Законодательная база и регулирование цифровых активов}

Управление цифровыми активами гражданина связано с рядом правовых аспектов, которые необходимо учитывать при разработке систем управления. В первую очередь, это вопросы собственности и прав доступа к цифровым активам. В различных юрисдикциях существуют разные подходы к определению прав собственности на цифровые активы, что может влиять на способы их управления и передачи.

В российской Федерации правовой статус цифровых активов не имеет чёткого правового определения, однако попытки регулирования в этой области предпринимаются. Например, уже действует закон о цифровых финансовых активах, который регулирует обращение криптовалют и токенов, обозначая их правовой статус, как имущественные права~\cite{federalLaw259FZ_2020}. Но в законе рассматривается только узкий сегмент финансовых цифровых активов и игнорируется широкий спектр других следов цифровой жизни граждан~\cite{Orlov2024Moluch}, таких как учетные записи в социальных сетях, электронная почта и т.д.

Кроме того, важным аспектом является защита персональных данных и конфиденциальности гражданина. Законодательства многих стран предусматривают меры по защите цифровых прав граждан, включая регулирование сбора, хранения и использования персональных данных, а, например, тайна переписки и право на частную жизнь закреплены в конституции Российской Федерации~\cite{constitutionRu}.

В мировой практике также существуют примеры регулирования цифровых активов. Например, в Европейском союзе действует Общий регламент по защите данных (GDPR), который устанавливает строгие правила обработки персональных данных и предоставляет гражданам контроль над своими цифровыми правами. Согласно этому регламенту, страны члены ЕС могут самостоятельно определять правила и порядок управления цифровыми активами в случае смерти владельца.

Таким правом воспользовались, например, во Франции, где в 2016 году был принят закон, позволяющий наследникам умершего получать доступ к его цифровым активам, включая учетные записи в социальных сетях и электронную почту.

Помимо этого, в ряде стран существуют специальные законы, регулирующие наследование цифровых активов. Например, в США действует Revised Uniform Fiduciary Access to Digital Assets Act, RUFADAA, который устанавливает порядок доступа и управления цифровыми активами после смерти владельца. Он предусматривает иерархию доступа к цифровым активам, основанную на воле владельца. Интересным является то, что согласно этому документу, самый высокий приоритет имеет воля владельца, указанная в соответствующем инструменте онлайн-сервиса, например в настройках аккаунта, если такая возможность предусмотрена~\cite{RUFADAA2015} и только во вторую очередь учитывается завещание или иные правовые документы.

\subsection{Проблемы и вызовы в правовом регулировании цифровых активов}

Несмотря на существующие законодательные инициативы, правовое регулирование цифровых активов сталкивается с рядом проблем и вызовов. Во-первых, это отсутствие единого международного стандарта, что приводит к разногласиям между юрисдикциями и усложняет управление цифровыми активами для граждан, имеющих активы в разных странах. Во-вторых, быстрое развитие технологий и появление новых типов цифровых активов создают сложности для законодательных органов, которые не успевают адаптировать свои нормы к новым реалиям. В-третьих, вопросы конфиденциальности и безопасности данных становятся все более актуальными, особенно в контексте управления цифровыми активами, что требует разработки эффективных механизмов защиты прав граждан.

В работе Орлова А.В. выделяются следующие ключевые проблемы правового регулирования цифровых активов граждан: неопределённость правового статуса цифровых активов, сложности в определении наследников и порядок передачи цифровых активов, а также вопросы конфиденциальности и безопасности данных~\cite{Orlov2024Moluch}.
Таким образом, правовые аспекты управления цифровыми активами граждан требуют дальнейшего развития и совершенствования законодательной базы, а также разработки международных стандартов и механизмов защиты прав граждан в цифровой среде.

